\section{Project Evolution}

\subsection{Phase 1: an educational backtesting scaffold}

The project began with BTCUSDT hourly OHLCV bars. Four simple strategy families
were implemented behind a common interface: momentum, mean reversion, moving
average crossover, and Bollinger bands. The surrounding framework separated
configuration, data loading, portfolio accounting, transaction costs, metrics,
plotting, parameter sweeps, and walk-forward orchestration.

That phase produced valuable engineering lessons: fees could erase a gross
edge, stateful strategy instances had to be recreated between runs, expanding
window selection changed parameter rankings, and orchestration code required as
much testing as individual components. Those lessons directly informed the L2
design.

The historical Phase 1 performance table is \emph{not} treated as rigorous
strategy evidence in this report. A later publication audit identified four
limitations:

\begin{enumerate}
  \item strategies observe the current bar close and execute at that same close,
  \item hourly equity observations were annualized with 365 periods per year,
  \item the displayed ``OOS Sharpe'' is a mean of monthly Sharpe estimates rather
        than the Sharpe of one stitched OOS return series, and
  \item the historical optimizer source was overwritten by the plotting pass.
        A runner has since been reconstructed from the logged grids, but the
        original performance artifacts are not promoted to current evidence.
\end{enumerate}

The current metrics helper now infers the observation frequency from its
timestamp index, and a regression test covers the hourly-versus-daily scaling.
The historical Phase 1 CSVs were not rewritten, so their published ratios
retain the original limitation.

These issues do not affect the independently implemented L2 replay. They do mean
that Phase 1 should be read as an engineering precursor and a record of learning,
not as evidence that any bar strategy generalizes. This distinction is now made
explicit rather than allowing precise-looking ratios to imply a stronger study.

\subsection{Why bars were insufficient}

Market making depends on phenomena that bars remove:

\begin{itemize}
  \item whether a limit order was passive when it reached the exchange,
  \item volume ahead at a price and how trades advance queue position,
  \item partial fills and residual inventory,
  \item entry and cancellation latency,
  \item event ordering when book changes and trades share a timestamp, and
  \item whether a fill is followed by a favorable or adverse price move.
\end{itemize}

The project therefore shifted from optimizing directional bar signals to
reconstructing a deterministic L2 event path and treating execution as part of
the hypothesis.

\subsection{Phase 2: replay, execution, and a frozen study}

Phase 2 added data recorders, snapshot-aware parsers, an L2 book, depth/trade
merging, execution simulation, passive quoting strategies, and markout/P\&L
analysis. Exploratory work selected a microprice reference, a 2 USDT
half-spread, and a five-second requote interval. Six anchor windows were later
retained inside a larger 24-window development panel. Because those anchors
were involved in exploration, the final panel is broader and pre-selected but
not a wholly untouched confirmation sample. The report preserves that fact.

The expanded study was organized as:

\begin{description}[leftmargin=5.4em,style=nextline]
  \item[Phase A] remeasure the fixed passive baseline at queue-credit endpoints,
  \item[Phase B] test whether OFI satisfies a population and fill-conditioned
        premise gate before running a strategy that uses it, and
  \item[Phase C] stress queue-cancellation credit and entry latency, then close
        the development arc if no candidate qualifies.
\end{description}

No candidate qualified. No candidate strategy was evaluated on the holdout.

\subsection{Phase 2.5: execution-model closure}

After the study was frozen, review of private-event timing found that the
historical engine used the next depth event as the activation point for an order
whose modeled arrival might have occurred earlier. Cancels were immediate, so
cancel/fill races were absent. Overlapping own orders also required a stronger
FIFO decomposition to guarantee that one recorded trade could not fill more
simulated volume than it contained.

The response was a provenance boundary:

\begin{itemize}
  \item historical artifacts were labeled \LegacyModel,
  \item the frozen result directory became write-protected to current scripts,
  \item the new implementation was named \CurrentModel,
  \item current results were assigned a distinct V3 namespace, and
  \item cross-model comparisons were restricted to descriptive sensitivity,
        not treated as estimates from the same execution process.
\end{itemize}

This is the current stage. The Python mechanics are tested, V3 development
economics remain pending.

\subsection{Timeline and decision discipline}

\begin{table}[H]
\centering
\caption{Condensed project timeline.}
\begin{tabularx}{\textwidth}{@{}L{0.15\textwidth}L{0.25\textwidth}Y@{}}
\toprule
Period & Milestone & Research consequence \\
\midrule
Jan--Apr 2026 & Bar framework and walk-forward exercise & Established reusable interfaces and exposed the importance of costs and state isolation. \\
Apr--May 2026 & L2 recorder, replay, execution, and exploratory baseline & Profit contaminated by accidental taker fills was rejected after post-only enforcement, a later positive one-window passive result failed broader-window evaluation. \\
May 2026 & Reconciliation, markouts, queue and same-ms audits & Shifted emphasis from headline P\&L to fill quality and model risk. \\
Jun 2026 & Frozen 24-window Phase A/B/C study & Baseline remained unattractive, OFI premise gate blocked the candidate, the holdout stayed strategy-sealed. \\
Jul 2026 & Private-event timing and own-FIFO correction & Froze legacy evidence and accepted \CurrentModel\ as the Python reference. \\
Aug 2026 & Public-release audit and technical report & Narrowed overstrong claims, demoted Phase 1 performance, and separated clone-level verification from full-data replay. \\
Next & V3 development-panel rerun & Determine whether current-model fills and economics preserve or change the historical conclusions. \\
\bottomrule
\end{tabularx}
\end{table}
